\section{模型假设}

为在保留系泊系统主要受力特征的基础上建立可求解的准静态模型，结合题目条件及项目统一口径，作如下假设。

\begin{enumerate}[label=\arabic*., leftmargin=2.8em]
  \item 将浮标视为轴线始终保持竖直的圆柱形刚体，仅考虑其竖直方向的吃水变化和水平方向的位移。风荷载等效作用于浮标水面以上迎风投影面积的形心处。由于题目未给出浮标重心位置、内部质量分布及稳性参数，忽略浮标在风荷载和系泊拉力作用下产生的倾斜。

  \item 题目未给出波高、波浪周期及波向等波浪参数，因此忽略波浪作用，将海面视为水平面，研究风速和流速作用下系泊系统的准静态平衡状态。该处理能够反映系统的平均姿态，但不描述极端海况下的瞬时倾角和锚链张力。

  \item 重物球采用均质实心铸钢球的假设，钢材密度取
  \(\rho_{\mathrm s}=7.85\times10^{3}\,\mathrm{kg/m^3}\)。设重物球质量为
  \(m_{\mathrm s}\)，则其体积和直径分别按
  \begin{equation}
    V_{\mathrm s}=\frac{m_{\mathrm s}}{\rho_{\mathrm s}},
    \qquad
    D_{\mathrm s}
    =\left(\frac{6m_{\mathrm s}}{\pi\rho_{\mathrm s}}\right)^{1/3}
    \label{eq:ball-volume-diameter}
  \end{equation}
  计算。问题一中 \(m_{\mathrm s}=1200\,\mathrm{kg}\)，由式
  \eqref{eq:ball-volume-diameter} 得
  \(V_{\mathrm s}=0.152866242\,\mathrm{m^3}\)，
  \(D_{\mathrm s}=0.663393159\,\mathrm{m}\)。

  \item 将连接重物球的悬挂绳简化为质量可忽略、不可伸长的柔性细绳，忽略悬挂绳自身所受的重力、浮力及水流阻力，认为其始终处于张紧状态，仅传递沿绳方向的拉力。本文采用准静态分析，不考虑重物球及悬挂绳的摆动、振动和瞬时动态效应。

  \item 锚链按均质钢材制造、且各链节均为游荡普通链环处理。沿用上述钢材密度；与项目给定的
  \(17.5\,\mathrm{mm}\) 链环钢材直径及 \(105\,\mathrm{mm}\) 单环长度相对应，链环宽度取
  \(61\,\mathrm{mm}\)。由附表中第 \(\kappa\) 型锚链的单位长度质量
  \(\mu_{\kappa}\) 和锚链长度 \(L\)，以
  \begin{equation}
    V_{\mathrm c}=\frac{\mu_{\kappa}L}{\rho_{\mathrm s}}
    \label{eq:chain-displaced-volume}
  \end{equation}
  计算锚链的等效排水体积；在此处理中不再将链环宽度重复计入浮力。
  各型号每米等效排水体积分别为
  \[
  \begin{aligned}
    \mathrm{I}:&\ 0.000407643\,\mathrm{m^3/m},&
    \mathrm{II}:&\ 0.000891720\,\mathrm{m^3/m},\\
    \mathrm{III}:&\ 0.001592357\,\mathrm{m^3/m},&
    \mathrm{IV}:&\ 0.002484076\,\mathrm{m^3/m},\\
    \mathrm{V}:&\ 0.003582166\,\mathrm{m^3/m}.
  \end{aligned}
  \]
  对问题一中的 II 型 \(22.05\,\mathrm{m}\) 锚链，其钢材体积为
  \(0.019662420\,\mathrm{m^3}\)，对应
  \(22.05/0.105=210\) 个链环。

  \item 锚的质量取题面给出的 \(600\,\mathrm{kg}\)，锚的抓地力按项目统一口径取
  \(132\,\mathrm{kN}\)。由于题目未给出海床摩擦系数、土壤类型及极限抗滑参数，本文不另行构造精确的锚滑移模型，而以题目给出的锚端锚链切线与海床夹角不超过
  \(16^\circ\) 作为抗拖移安全判据。

  \item 重力加速度统一取
  \(g=9.8\,\mathrm{m/s^2}\)。浮标、钢管、钢桶、重物球、锚链及锚的重力均按质量乘以该重力加速度计算。

  \item 从浮标向下将四节钢管依次编号为第 \(1,2,3,4\) 节。每节钢管视为不可伸长且不发生弯曲的刚性杆；相邻钢管之间的连接简化为无摩擦铰接，连接处允许相对转动，仅传递作用力而不传递弯矩。
\end{enumerate}

上述假设中，浮标和钢桶的姿态分别通过其轴线位置与竖直方向的夹角描述；锚链末端的安全约束采用其切线方向与水平海床的夹角描述。除特别说明外，长度、质量和力分别采用国际单位制中的米、千克和牛顿，角度在结果展示时统一换算为度。
